\chapter[Sermon 31. Here Is Opened the iii. and iiii. Seal, and Is Declared What the World Shall Suffer of Hunger and Pestilence.]{Here Is Opened the iii. and iiii. Seal, and Is Declared What the World Shall Suffer of Hunger and Pestilence.}
\label{sermon:031}
\markright{Sermon 31. Here Is Opened the iii. and iiii. Seal, and Is ...}

And when he opened the third seal, I heard the third beast say, “Come and see.” And I looked, and behold, a black horse; and he who sat on him had a pair of balances in his hand. And I heard a voice in the midst of the four beasts saying, “A measure of wheat for a penny, and three measures of barley for a penny; and injure not the oil and wine.”

And when he opened the fourth Seal, I heard the voice of the fourth living creature say, “Come and see.” And I looked, and behold a pale horse, and the name of the one who sat on him was Death, and Hell followed after him, and power was given to them over the fourth part of the earth, to kill with sword, and with famine, and with death of the beasts of the earth.

Christ, exalted above all things, and Lord of all in heaven and on earth, opens the seals of the divine book, that is to say, disposes and governs with great righteousness the ordinances and judgments of God. First, indeed, he gives a prosperous course to the preaching of the gospel, sending always faithful ministers preaching the Gospel of the kingdom of God, peace, and concord. But because evil men despise the Evangelical peace, they are certainly worthy to be afflicted with cruel wars. Therefore, the Lamb opens the second seal, and there rush out cruel wars, slaughters, seditions, and robberies.

\index{Repentance}But before the third Seal is opened, the third beast, resembling the face of a man, exhorts us to take most diligent heed: that when we see these things come to pass which are spoken of here, we should consider from whence they come, and for what causes they are sent, and that they may be turned away by true repentance. Some refer these things absolutely to chance and fortune, others to natural causes, with no respect had at all for God and the divine operation: whereas we know that God uses natural causes according to his good will and pleasure. Let us watch therefore, look and consider, and know that the righteous God works all things for the salvation of the elect, and the overthrow of his enemies. That black horse with its rider, showing a balance in his hand, signifies the unfortunate or sorrowful time of scarcity, famine, and penury of all things.

For it is a worthy and fitting punishment that those who do not esteem the bread of life, nor have any consideration of the food of souls, but both reject it themselves, and by their tyrannical proclamations bring it to pass that it is not received by others, and finally who, for the bread of life, despoil the godly of their goods, and most wickedly waste them in all kinds of riot, should be driven to necessities at excessive prices: yea, and cannot find necessities, but should pine for hunger. We know that the color black is used in mourning and grief; and that when flesh and blood are consumed for want of food, the skin grows black and ill-favored: and therefore this horse is black.

\index{John, Saint}\index{Erasmus of Rotterdam}\index{Arethas of Caesarea}The rider of this horse holds in his hand a balance, [?a token] with two scales hanging at either end of the beam, which we call a pair of weights. Arethas says that a balance is a token of right and equity. For you have sat, says David, upon your throne which judges righteousness; therefore a balance is the judgment of the just judgment of God. Arethas has not alleged these things amiss, but we ought rather to prefer the exposition of Saint John himself. For a voice is heard from the midst of the beasts, which expounds to us the balance. For it sounds, a measure of wheat for a penny, and three measures of barley for a penny. And this measure, called Choinix, signifies a diet or daily meat, as Erasmus has in his proverb, “Sit not upon your measure.” The same in his annotations upon this place: Choinix, says he, is a measure of wheat, or other breadcorn, which is sufficient for one day’s meat. Budaeus thinks that it weighs four pounds, Pollux three. The word therefore signifies that a very little meat shall cost a great price, and yet not be gotten for money, which happens in the time of famine. What the Roman penny is worth Budaeus shows; we understand by it plainly a great price. Therefore two things are signified: scarcity or dearth of corn, and famine. Scarcity raises the price beyond reason. Famine has nothing to buy, though he has never so much money lying by him, but hungers, wants, pines, and at the last miserably consumes to naught, wherein verily dearth and famine do differ. The Germans discern them by several words, calling dearth scarcity, and famine hunger. Yet they are for the most part indivisible.

And we read in the old story of the Bible, that for the contempt of the preaching of God’s law, and the bringing in of a strange kind of worshiping God, the Israelites in the times of Elijah and Elisha were most grievously punished with hunger and poverty. These things are plentifully declared in the third book of Kings, chapters 17 and 18. Also in the fourth of Kings, chapters 6 and 7. Moreover, in the time of Emperor Claudius, while the Apostles preached the Gospel faithfully, and the Jews and Gentiles stoutly repulsed it, famine most grievously afflicted the Roman Empire: which thing Luke rehearses in the Acts of the Apostles, chapter 11. These things were indeed done before this revelation was exhibited to John. Since that time, historians recite sundry and innumerable famines, dearths, and poverties in diverse countries, sent of God for contempt of the truth. Nauclerus mentions a famine in the year of our Lord 1339, wherein mothers also devoured their own children. What has chanced in our memory in those wars of Milan and elsewhere, it is no need to rehearse. They are yet fresh in memory, and written in the histories of Galeazzo Capella. We felt some part hereof also in the year of our Lord 1529 and the years following. The just Lord punishes, and more will punish the great ingratitude and contempt of his godly word, as he did in the destruction of Jerusalem. Would God it would please the world, so blind, through repentance to convert unto God when he punishes, and with free and willing minds embrace the word of truth: for so should there be more felicity and less misery.

However, for comfort at the end of this seal, it is added, “See that you hurt not the oil and the wine.” He names the kinds most necessary for human use, and means that God mercifully reserves some things that are chiefly necessary for humankind, especially for the sake of the elect, so that all should not perish and suffer generally. By this, we understand that the Lord does not forget his mercy, even in the midst of affliction and plagues that he sends. Thus, in times past, intending to punish Egypt and other nations with famine, he sent Joseph beforehand, by whom he might preserve the house of Jacob and other people innumerable. You see herein most clearly that it is from God that sometimes the grain is blasted, and the vines and olives perish; and that it is from him that the grain increases, and wine also. So he has also previously protested in the law, Leviticus 26, and Deuteronomy 28.

We have now come to the fourth Seal, at the opening of which, and to observe its effect, we are stirred up by the Eagle, the fourth beast. Of whom we have spoken before on one or two occasions. And the pale Horse comes forth, in Greek \textit{chloros}, a color resembling withered grass and herbs. Solomon in the twelfth chapter of Ecclesiastes calls the color appearing in dead bodies and their faces a golden liquid. All poets call death pale. And the rider is indeed expressly called Death. We understand the course of plague and all diseases, and even death itself; and Hell follows, that is to say, a pit or a grave. For \textit{Sheol} in Hebrew signifies a pit or a grave. But if you will needs understand it of the place of the damned, doubtless they are carried headlong into Hell, so many as here, consumed by sickness, die without faith and repentance. Therefore Hell follows death rightly. But if you would rather understand Hell as a grave, it signifies that it will be full of corpses and sepulchers.

\index{Rome}And indeed plagues and pestilences most mortal have severely afflicted the Roman Empire. Orosius testifies in his seventh book in the Acts of L. Aurelius Verus and Decius, emperors, the most cruel persecutors of our faith. Evagrius in chapter 29 of book 4 of Ecclesiastical History tells of a marvelous plague that lasted about 50 years. And all men know with what a pestilence and sudden death Italy was wasted in the time of Emperor Maurice and Bishop Gregory of Rome. Time would fail me if I were to recount from histories all the plagues and calamities of all times. What is done at this day, and has been done in our memory, you yourselves know best. New diseases have arisen, whose names were never known to our elders. With these evils and calamities God wastes the world, and has always done so, intending that by plagues he might call us again to repentance. Thus we must always judge of calamities. If anyone judges otherwise, they are not reformed, therefore they are punished here and after this will burn in perpetual torments.

\index{Ezekiel (Book of)}To these moreover is added another thing also, and power was given them, and so forth. For when people will not reform with simple calamities, the evils or plagues of God are doubled. They are listed in the same order and number as the prophets, Jeremiah in chapter 15 and Ezekiel in chapter 14. For they are these: sword, famine, death or pestilence, and wild beasts; so they are recited in the Law also. With these, as if sent from the four parts of the world, God, most righteous, executes his judgments.

Let us observe this chiefly: that power is given by God to kill, and that over a fourth part of the earth. For we learn that God alone quickens and slays, and that he works this justly by his instruments, finally that all his things are numbered and done in order. Whereupon he pours out his fury upon a third part of the world. For he knows whom he should punish and whom he should nurture tenderly.

\index{Eusebius of Caesarea}Certainly, accounts attest how in dire circumstances, when all things are brought to an extremity of misfortune, God has unleashed sword, pestilence, famine, and wild beasts to afflict people. And Arethas appropriately cites the words of his predecessor, Saint Andrew, Bishop of Caesarea, from the \textit{Ecclesiastical History} of Eusebius, book 9, chapter 8. Indeed, within the last fifty years, historians have reported many such occurrences, and we have witnessed some ourselves.

Therefore, if we desire to be free from such great evils, let us serve God in truth and value his word, which he has sent to heal us. It is reasonable that those who reject sound doctrine should be afflicted with various diseases of soul and body, and so forth.

\index{Augustine, Saint}You will say, but these evils invade even the best of things. And indeed they do. Why God permits this, Augustine explains at length in the first book of \textit{The City of God}. Certainly, for the godly, all things work for the best. The thieves suffered the same death on the cross that Christ did, and he as they: but the consideration of them is very different. The apostles and countless martyrs die by the sword, just as soldiers do in war, but with unlike circumstances. The godly become partakers of the suffering of the Son of God. The ungodly are punished for their wickedness, and their suffering is without glory; rather, this is the beginning, unless they acknowledge him who strikes them, of everlasting torments. The Lord preserve us from evil.
